% GENERATED FILE. Edit canonical lessons, then run npm run generate:books.
% canonical-source-hash: fnv1a64:cf4dec85aa98de03
% canonical-lessons: PA-C354-kahani, PA-C354-sunauna, PA-C354-kissa, PA-C354-rani, PA-C354-jadu

\chapter{Once Upon a Time}
\label{ch:pa-a2-once-upon-a-time}
\hlchaptermodality{354}

\begin{chapteropening}
\textbf{By the end of this chapter:} \emph{I can say which stories I like and who tells them.}

Complete the last lesson of chapter 354: I can say which stories I like and who tells them.
\end{chapteropening}

\section[\texorpdfstring{\pa{ਕਹਾਣੀ} (kahāṇī)}{ਕਹਾਣੀ (kahāṇī)}]{\pa{ਕਹਾਣੀ} (kahāṇī) — a story}

\label{lesson:PA-C354-kahani}



\begin{quote}
Before the new one: say the Punjabi for to make something fly, then the Punjabi for to play with a child.
\end{quote}

\subsection*{You'll want to know: \pa{ਕਹਾਣੀ}}
\textbf{\pa{ਕਹਾਣੀ}} — \emph{kahāṇī} — \textquotedblleft{}a story\textquotedblright{}.

\begin{quote}
\pa{ਮੈਨੂੰ} \pa{ਪੁਰਾਣੀਆਂ} \pa{ਕਹਾਣੀਆਂ} \pa{ਪਸੰਦ} \pa{ਹਨ।}
\end{quote}

\begin{quote}
\emph{mainū̃ purāṇīāṁ kahāṇīāṁ pasand han.}
\end{quote}

\begin{quote}
I like old stories.
\end{quote}

A feminine noun.

\begin{cousinweb}
\textbf{\pa{ਕਹਾਣੀ}} (\emph{kahāṇī}) is inherited from Sanskrit \emph{kathānaka}, \textquotedblleft{}a little tale\textquotedblright{}; CDIAL lists Punjabi \emph{kahāṇī} under it (CDIAL 2705).
\end{cousinweb}

\begin{grammarlens}[title={What you've built}]
One more word for this chapter.
\end{grammarlens}

\subsection*{Your turn}
\begin{itemize}
\raggedright
  \item \emph{Say it:} \emph{kahāṇī}
  \item \emph{Say it:} \emph{kahāṇī}, once more
  \item \emph{Say it:} \emph{mainū̃ purāṇīāṁ kahāṇīāṁ pasand han.} — I like old stories.
  \item \emph{Recall:} say the Punjabi for to make something fly, then the Punjabi for to play with a child, then say \emph{kahāṇī} again
\end{itemize}

\begin{tcolorbox}[breakable,colback=teal!4,colframe=teal!35!black,title={Before you move on}]
What does \pa{ਕਹਾਣੀ} mean? (\textquotedblleft{}A story\textquotedblright{}.) Say it once more.
\end{tcolorbox}
\section[\texorpdfstring{\pa{ਸੁਣਾਉਣਾ} (suṇāuṇā)}{ਸੁਣਾਉਣਾ (suṇāuṇā)}]{\pa{ਸੁਣਾਉਣਾ} (suṇāuṇā) — to tell a story; to recite}

\label{lesson:PA-C354-sunauna}



\begin{quote}
Before the new one: say the Punjabi for to play with a child, then the Punjabi for a story.
\end{quote}

\subsection*{You'll want to know: \pa{ਸੁਣਾਉਣਾ}}
\textbf{\pa{ਸੁਣਾਉਣਾ}} — \emph{suṇāuṇā} — \textquotedblleft{}to tell a story; to recite\textquotedblright{}.

\begin{quote}
\pa{ਦਾਦੀ} \pa{ਰਾਤ} \pa{ਨੂੰ} \pa{ਕਹਾਣੀ} \pa{ਸੁਣਾਉਂਦੀ} \pa{ਹੈ।}
\end{quote}

\begin{quote}
\emph{dādī rāt nū̃ kahāṇī suṇāundī hai.}
\end{quote}

\begin{quote}
Grandmother tells a story at night.
\end{quote}

Past (I …): \pa{ਮੈਂ} \pa{ਇੱਕ} \pa{ਕਹਾਣੀ} \pa{ਸੁਣਾਈ}, \emph{maiṁ ikk kahāṇī suṇāī}.

\begin{cousinweb}
\textbf{\pa{ਸੁਣਾਉਣਾ}} (\emph{suṇāuṇā}) is the causative partner of \emph{suṇnā}, \textquotedblleft{}to hear\textquotedblright{}, which goes back to Sanskrit \emph{śṛṇoti}, \textquotedblleft{}hears\textquotedblright{}; CDIAL gives no Punjabi form, but lists Lahnda \emph{suṇāvaṇ}, \textquotedblleft{}informs\textquotedblright{}, under it (CDIAL 12598).
\end{cousinweb}

\begin{grammarlens}[title={What you've built}]
One more word for this chapter.
\end{grammarlens}

\subsection*{Your turn}
\begin{itemize}
\raggedright
  \item \emph{Say it:} \emph{suṇāuṇā}
  \item \emph{Say it:} \emph{suṇāuṇā}, once more
  \item \emph{Say it:} \emph{dādī rāt nū̃ kahāṇī suṇāundī hai.} — Grandmother tells a story at night.
  \item \emph{Recall:} say the Punjabi for to play with a child, then the Punjabi for a story, then say \emph{suṇāuṇā} again
  \item \emph{Write it:} copy \textbf{\pa{ਸੁਣਾਉਣਾ}} once, checking the whole word
\end{itemize}

\begin{tcolorbox}[breakable,colback=teal!4,colframe=teal!35!black,title={Before you move on}]
What does \pa{ਸੁਣਾਉਣਾ} mean? (\textquotedblleft{}To tell a story; to recite\textquotedblright{}.) Say it once more.
\end{tcolorbox}
\section[\texorpdfstring{\pa{ਕਿੱਸਾ} (kissā)}{ਕਿੱਸਾ (kissā)}]{\pa{ਕਿੱਸਾ} (kissā) — a tale; a long story}

\label{lesson:PA-C354-kissa}



\begin{quote}
Before the new one: say the Punjabi for a story, then the Punjabi for to tell a story.
\end{quote}

\subsection*{You'll want to know: \pa{ਕਿੱਸਾ}}
\textbf{\pa{ਕਿੱਸਾ}} — \emph{kissā} — \textquotedblleft{}a tale; a long story\textquotedblright{}.

\begin{quote}
\pa{ਦਾਦਾ} \pa{ਇੱਕ} \pa{ਲੰਮਾ} \pa{ਕਿੱਸਾ} \pa{ਸੁਣਾਉਂਦਾ} \pa{ਹੈ।}
\end{quote}

\begin{quote}
\emph{dādā ikk lammā kissā suṇāundā hai.}
\end{quote}

\begin{quote}
Grandfather tells a long tale.
\end{quote}

A masculine noun.

\begin{cousinweb}
\textbf{\pa{ਕਿੱਸਾ}} (\emph{kissā}) is a loan, usually traced through Persian to Arabic \emph{qiṣṣa}, \textquotedblleft{}story\textquotedblright{}.
\end{cousinweb}

\begin{grammarlens}[title={What you've built}]
One more word for this chapter.
\end{grammarlens}

\subsection*{Your turn}
\begin{itemize}
\raggedright
  \item \emph{Say it:} \emph{kissā}
  \item \emph{Say it:} \emph{kissā}, once more
  \item \emph{Say it:} \emph{dādā ikk lammā kissā suṇāundā hai.} — Grandfather tells a long tale.
  \item \emph{Recall:} say the Punjabi for a story, then the Punjabi for to tell a story, then say \emph{kissā} again
  \item \emph{Write it:} copy \textbf{\pa{ਕਿੱਸਾ}} once, checking the whole word
\end{itemize}

\begin{tcolorbox}[breakable,colback=teal!4,colframe=teal!35!black,title={Before you move on}]
What does \pa{ਕਿੱਸਾ} mean? (\textquotedblleft{}A tale; a long story\textquotedblright{}.) Say it once more.
\end{tcolorbox}
\section[\texorpdfstring{\pa{ਰਾਣੀ} (rāṇī)}{ਰਾਣੀ (rāṇī)}]{\pa{ਰਾਣੀ} (rāṇī) — a queen}

\label{lesson:PA-C354-rani}



\begin{quote}
Before the new one: say the Punjabi for to tell a story, then the Punjabi for a tale.
\end{quote}

\subsection*{You'll want to know: \pa{ਰਾਣੀ}}
\textbf{\pa{ਰਾਣੀ}} — \emph{rāṇī} — \textquotedblleft{}a queen\textquotedblright{}.

\begin{quote}
\pa{ਕਹਾਣੀ} \pa{ਵਿੱਚ} \pa{ਇੱਕ} \pa{ਸਿਆਣੀ} \pa{ਰਾਣੀ} \pa{ਹੈ।}
\end{quote}

\begin{quote}
\emph{kahāṇī vich ikk siāṇī rāṇī hai.}
\end{quote}

\begin{quote}
There is a wise queen in the story.
\end{quote}

A feminine noun.

\begin{cousinweb}
\textbf{\pa{ਰਾਣੀ}} (\emph{rāṇī}) is inherited from Sanskrit \emph{rājñī}, \textquotedblleft{}queen\textquotedblright{}; CDIAL lists Punjabi \emph{rāṇī}, \textquotedblleft{}queen\textquotedblright{}, under it (CDIAL 10692).
\end{cousinweb}

\begin{grammarlens}[title={What you've built}]
One more word for this chapter.
\end{grammarlens}

\subsection*{Your turn}
\begin{itemize}
\raggedright
  \item \emph{Say it:} \emph{rāṇī}
  \item \emph{Say it:} \emph{rāṇī}, once more
  \item \emph{Say it:} \emph{kahāṇī vich ikk siāṇī rāṇī hai.} — There is a wise queen in the story.
  \item \emph{Recall:} say the Punjabi for to tell a story, then the Punjabi for a tale, then say \emph{rāṇī} again
  \item \emph{Write it:} copy \textbf{\pa{ਰਾਣੀ}} once, checking the whole word
\end{itemize}

\begin{tcolorbox}[breakable,colback=teal!4,colframe=teal!35!black,title={Before you move on}]
What does \pa{ਰਾਣੀ} mean? (\textquotedblleft{}A queen\textquotedblright{}.) Say it once more.
\end{tcolorbox}
\section[\texorpdfstring{\pa{ਜਾਦੂ} (jādū)}{ਜਾਦੂ (jādū)}]{\pa{ਜਾਦੂ} (jādū) — magic}

\label{lesson:PA-C354-jadu}



\begin{quote}
Before the new one: say the Punjabi for a tale, then the Punjabi for a queen.
\end{quote}

\subsection*{You'll want to know: \pa{ਜਾਦੂ}}
\textbf{\pa{ਜਾਦੂ}} — \emph{jādū} — \textquotedblleft{}magic\textquotedblright{}.

\begin{quote}
\pa{ਉਸ} \pa{ਕਹਾਣੀ} \pa{ਵਿੱਚ} \pa{ਬਹੁਤ} \pa{ਜਾਦੂ} \pa{ਹੈ।}
\end{quote}

\begin{quote}
\emph{us kahāṇī vich bahut jādū hai.}
\end{quote}

\begin{quote}
There is a lot of magic in that story.
\end{quote}

A masculine noun.

\begin{cousinweb}
\textbf{\pa{ਜਾਦੂ}} (\emph{jādū}) is a loan, usually traced to Persian \emph{jādū}, \textquotedblleft{}magic\textquotedblright{}.
\end{cousinweb}

\begin{grammarlens}[title={What you've built}]
One more word for this chapter.
\end{grammarlens}

\subsection*{Your turn}
\begin{itemize}
\raggedright
  \item \emph{Say it:} \emph{jādū}
  \item \emph{Say it:} \emph{jādū}, once more
  \item \emph{Say it:} \emph{us kahāṇī vich bahut jādū hai.} — There is a lot of magic in that story.
  \item \emph{Recall:} say the Punjabi for a tale, then the Punjabi for a queen, then say \emph{jādū} again
  \item \emph{Write it:} copy \textbf{\pa{ਜਾਦੂ}} once, checking the whole word
\end{itemize}

\begin{tcolorbox}[breakable,colback=teal!4,colframe=teal!35!black,title={Before you move on}]
What does \pa{ਜਾਦੂ} mean? (\textquotedblleft{}Magic\textquotedblright{}.) Say it once more.
\end{tcolorbox}
